\icode{E-EXE-FI-L3-261001-001}
某同学根据图甲所示的欧姆表原理图，利用微安表（满偏电流\,\qty{300}{\mu A}，内阻\,\qty{100}{\Omega}）、滑动变阻器\(R\)（最大阻值为\,\qty{10}{k\Omega}）和一节干电池，将微安表改装成欧姆表。

	\begin{figure}[htbp]
		\centering
		\begin{tikzpicture}
			\draw[thick] (-.1,3.9)--(-.1,4.1) (.1,3.8)--(.1,4.2);
			\draw[blue] (.1,4)--(1.5,4)--(1.5,1.5)--(1,1.5)--(1,1) (-.1,4)--(-1.5,4)--(-1.5,1.5)--(-1,1.5)--(-1,1);
			\draw[->,blue] (1.5,2)--(2,2)--(2,3)--(1.6,3);
			\fill (1.5,2) circle(.05);
			\draw[red,very thick] (-1,1)--(-1,0)node[below]{\(A\)};
			\draw[very thick] (1,1)--(1,0)node[below]{\(B\)};
			\filldraw[fill=white] (-1.5,3)node{\(\mu \mathrm{A}\)} circle(.3) (1.4,2.6) rectangle(1.6,3.4) (-1,1) circle(.05) (1,1)circle(.05);
			\node at(1.4,3)[left]{\(R\)};
			\node at(0,3.8)[below]{\(E,r\)};
			\draw[dotted] (-2,1.25) rectangle(2.5,4.45);
			\node at(0,-.5)[below]{\footnotesize 图甲};
		\end{tikzpicture}
		\begin{tikzpicture}[scale=1.2]
			\draw (60:4) arc(60:120:4);
			\foreach \x in {60,62,...,120}
				\draw (\x:4)--(\x:4.1);
			\foreach \x in {70,90,110}
				\draw[blue] (\x:4)--(\x:4.2);
			\foreach \x/\n in {60/300,80/200,100/100,120/0}
				\draw[red] (\x:4)--(\x:4.3)node[above]{\footnotesize \n};
			\node at(0,3) {\small\(\mathrm{\mu A}\)};
			\draw (80:2.5)--(80:4);
			\node at(0,1.5)[below]{\footnotesize 图乙};
		\end{tikzpicture}
	\end{figure}

	(1)将两表笔短接，调节\(R\)使微安表指针在\lunl{4em}\,\qty{}{\mu A}处。

	(2)当两表笔之间接入阻值为\,\qty{2.4}{k\Omega}的定值电阻时，微安表指针在如图乙所示的位置，则其读数为\lunl{4em}\,\qty{}{\mu A}，电池的电动势为\lunl{4em}\,\qty{}{V}，欧姆表的内阻为\lunl{4em}\,\qty{}{k\Omega}。

	(3)将微安表上的“\(\mu \mathrm{A}\)”处标明“\(\Omega\)”，“300”位置处标明“0”，“100”位置标明“\lunl{4em}”，并在其他位置标明相应的电阻刻度值，这样就把微安表改装成了欧姆表。

	(4)经过一段时间之后，电池的电动势降低，内阻增大，则重新欧姆调零之后，测得的电阻阻值将\lunl{4em}（填“偏大”“不变”或“偏小”）。